\printedpage{59}
% head: A. Grothendieck : Dérivateurs, Ch. XV, Édition des Universités de Montpellier II et Paris VII

\begin{itemize}[label={},leftmargin=4.5em]
\item \textit{and the localized functors}
\[
\begin{tikzcd}
\mathrm{Hot}_{B,W} \arrow[r, shift left, "\overline{i_B}"] & \mathrm{Hot}_W \arrow[l, shift left, "\overline{j_B}"]
\end{tikzcd}
\]
\textit{are equivalences of categories. Moreover, these equivalences of categories are quasi-inverse to one another, the adjunction morphisms being deduced from the adjunction morphisms for the pair of adjoint functors $(i_B, j_B)$, these adjunction morphisms being moreover in $W$ resp. in $W_B$.}\footnote[38]{\textbf{N.B.} Condition $(d)$, seemingly very strong, is equivalent to the seemingly more innocuous condition:\par
$d'$) the functor $j_A : \mathrm{Cat} \longrightarrow A^\wedge$ transforms elements of $W$ into $W_A$-aspheric arrows;\par\noindent
and it even suffices to require it for $C \longrightarrow e$ in Cat, where $C$ is an object of Cat having a final object.\par\noindent
\texttt{[Condition ($d'$) above is too strong and is not equivalent to condition ($d$); on the other hand the weaker condition considered next is indeed equivalent to ($d$)].}}
\end{itemize}

\msp{page 107}

\begin{enumerate}[leftmargin=2.5em]
\item[$2^\circ$)] \textit{If the equivalent conditions of $1^\circ$ are satisfied, then the triple}
\[
(W_A\,,\ C = \mathrm{Mon}_{A^\wedge}\,,\ F = (C \cap W_A)_*)
\]
\textit{is a (closed) left proper Quillen triple.}
\item[$3^\circ$)] \textit{If the equivalent conditions of $1^\circ$ are satisfied, and if moreover $A$ is $W$-aspheric, then (and only then) $A$ itself is a} test category \textit{for $W$, which amounts to saying that the conditions on the category $B$ in} (d) \textit{are satisfied when $B$ is} either \textit{$A$, or a localized category $A/a$.}
\end{enumerate}

\textbf{Corollary 8.5.} \textit{The equivalent conditions of $1^\circ$ and of $3^\circ$ are satisfied in each of the following cases (and this for any $W$, satisfying the stated hypotheses).}
\begin{enumerate}[leftmargin=2.5em]
\item[a)] $A = \Simp$, $\Simp^{\sim}$, $\Cub$, $\Cub^{\sim}$ \textit{(categories of standard combinatorial simplices, ordered or not, and of standard combinatorial cubes, ordered or not \ldots)}\footnote[39]{$\Cub^{\sim}$ is essentially the ``universal'' category $A$, stable under finite products, and containing a ``separating interval'' $(I,\delta_0,\delta_1)$ as in $(b)$.}.
\item[b)] \textit{$A$ is stable under finite products}\footnote[40]{In fact, it suffices that $I \longrightarrow e$ be squarable, \textit{i.e.} that $A$ have a final object, and that the products $I \times a$ ($a \in \mathrm{Ob}\,A$) exist in $A$.}\textit{, and contains an object $I$ such that $I$ has two sections $\delta_0$, $\delta_1$ over the final object, with $\delta_0$, $\delta_1$ disjoint,} i.e. $\mathrm{Ker}(\delta_0,\delta_1) = \emptyset_{A^\wedge}$ \textit{(empty subpresheaf of the final presheaf).}
\end{enumerate}

\msp{page 108}

Case $(b)$ of the corollary falls under condition $(c)$ in $1^\circ$. (Moreover, $A$ having a final object, is $W$-aspheric, so we are under the conditions of $3^\circ$.) The cases $\Simp$, $\Simp^{\sim}$, $\Cub$ (and $\Cub^{\sim}$, which in fact falls under case $(b)$) are treated in Pursuing Stacks \S34 (p. 47-50) \texttt{[see also \S36]}.

\textsc{Proof of theorem} 8.4. The essential part, technically, is of course $1^\circ$\footnote[41]{Part $2^\circ$ follows from it, thanks to corollary 8.3.}. We have the $\pm$ tautological implications
\[
(a) \ \Longrightarrow\ (b) \ \Longrightarrow\ (c)\ ,
\]
where $(a) \Longrightarrow (b)$ comes from the fact that $L$ is obviously an injective object of $A^\wedge$. (The Lawvere object of any topos is an injective object.) It remains to establish
\[
(c) \ \Longrightarrow\ (d) \ \Longrightarrow\ (a)\ .
\]
% foot: Version 5 novembre 2010

\printedpage{60}
% head: A. Grothendieck : Dérivateurs, Ch. XV, Édition des Universités de Montpellier II et Paris VII

That we have $(c) \Longrightarrow (d)$ (and $(c) \Longrightarrow (b)$) is established in Pursuing Stacks, see the summary of the $W$-test notions in Pursuing Stacks, \S65, p. 168-179; the relevant statement here is proposition 3, page 173.

It remains to establish $(d) \Longrightarrow (a)$ – but now we're on easy street! I note here the following lemma, whose place would rather be at the end of \S4 of the generalities on Quillen triples:

\msp{page 109}

\textbf{Lemma 8.6.} \textit{Let $\mathcal{M}$ be a category, $(C,TF)$ a Quillen precouple in $\mathcal{M}$, and $W \subset \mathrm{Fl}(\mathcal{M})$ a subset of} $\mathrm{Fl}(\mathcal{M})$ \textit{stable under direct factors. Equivalent conditions:}
\begin{enumerate}[leftmargin=2.5em]
\item[a)] \textit{$TF \subset W$.}
\item[b)] \textit{Every arrow $f$ in $\mathcal{M}$ factors as $f = pi$, with $i \in C$, $p \in W$.}
\item[b$'$)] \textit{As in} (b), \textit{with $f$ in $TF$.}
\end{enumerate}

\textsc{Proof}. The implications $(a) \Longrightarrow (b) \Longrightarrow (b')$ are tautological and do not appeal to the stability hypothesis on $W$. But for a given $f \in TF$, to see that it is in $W$, if $W$ is stable under direct factors, it suffices for that particular $f$ to factor as $pi$, with $i \in C$, $p \in W$. Indeed, we see at once that then $f$ is a direct factor of $p$ (since $i \in C$), and since $p \in W$, we have $i \in W$, \quad \textit{q.e.d.} % [sic] $i \in W$ -> $f \in W$, the conclusion to be proved concerns f

Let us return to 8.4, proof of $(d) \Longrightarrow (a)$.

\texttt{[There follows a crossed-out passage, at the bottom of page 109 and at the top of page 110.]}

\msp{page 110}

In the end, my plan of proof changes, for although $(d)$ looks very, very strong as a condition, as it stands I do not see, in the end, how to deduce $(a)$ from it. But in fact we can consider that the equivalence of $(b)$, $(c)$, $(d)$ is put here for the record, it is established in \textit{loc. cit.}, so I only have to concern myself with the implication $(c) \Longrightarrow (a)$, which will follow
\msp{page 111}
from lemma 8.6, which reduces us to proving that every arrow $f : X \longrightarrow Y$ in $A^\wedge$ factors as $pi$, with $i$ a monomorphism and $p \in W_A$. We use the mapping cone argument \texttt{[rather mapping cylinder]}, starting from
\[
\begin{tikzcd}[column sep=small]
e_A \amalg e_A \arrow[rr, hook, "{(\delta_0,\delta_1)}", "\text{mono}"'] \arrow[dr] & & I \arrow[dl, "(W_A)_{\mathrm{univ}}"] \\
& e_A\ , &
\end{tikzcd}
\]
whence for every $X$ in $A^\wedge$, taking the product with $X$,
\[
\begin{tikzcd}[column sep=small]
X \amalg X \arrow[rr, hook, "{(\delta_0^X,\delta_1^X)}", "\text{mono}"'] \arrow[dr, "{(\mathrm{id},\mathrm{id})}"] & & X \times I \arrow[dl, "\mathrm{pr}_1^X \in W_A"] \\
& X\ , &
\end{tikzcd}
\]
% foot: Version 5 novembre 2010
\printedpage{61}
% head: A. Grothendieck : Dérivateurs, Ch. XV, Édition des Universités de Montpellier II et Paris VII
and we construct the mapping cone \texttt{[rather mapping cylinder]} of $f$ by the cocartesian diagram
\[
\begin{tikzcd}[column sep=4em,row sep=4em]
 & X \arrow[r, "f"] \arrow[d, hook, shift right, "\delta_0^X"'] \arrow[dr, phantom, "\text{cocart.}"] & Y \arrow[d, hook, shift right, "\beta"'] \\
X \arrow[r, dotted, "\delta_1^X"] \arrow[rr, dotted, bend right, "i_f"'] & X \times I \arrow[r, "f'"] \arrow[u, dotted, shift right, "\mathrm{pr}_1^X"'] & Z(f)\ , \arrow[u, dotted, shift right, "p_f"']
\end{tikzcd}
\qquad
\begin{array}{c}
p_f \text{ deduced from} \\
\mathrm{pr}_1^X : X \times I \to X
\end{array}
\tag*{(8.6.1)}
\]
the desired factorization of $f$ is
\[
f \ = \ p_f\, i_f\ , \qquad \text{where}\quad i_f = f'\delta_1^X\ .
\]
Here, since $\mathrm{pr}_1^X \in W_A$, we conclude that
\msp{page 112}
$\delta_0^X$, $\delta_1^X$ (which are sections of it) are in $W_A$, since $\delta_0^X$ is a monomorphism, we thus have $\delta_0^X \in TC$, hence ($TC$ being stable under cobase change) $\beta \in W_A$, hence finally $p_f$ (which is a retraction of $Z(f)$ onto $Y$) is also in $W_A$. It remains to see that $i_f$ is a monomorphism, which is seen from the more complete diagram (\textit{cf.} XII, p. 230) that I have in front of me: $i_f$ appears as a composite of two monomorphisms $X \longrightarrow Y \amalg X$ and $Y \amalg X \longrightarrow Z(f)$ \ldots{} Theorem 8.4 is proved.

I will all the same isolate the everlasting formal argument (which I have used mostly in dual form) in a

\textbf{Lemma 8.7.} \textit{Let $\mathcal{M}$ be a category, $W$ a mildly saturated localizer in $\mathcal{M}$, $C$ a subset stable under composition and under cobase change. We propose, for $X$ fixed in $\mathcal{M}$, to find a functorial factorization of every arrow $f : X \longrightarrow Y$} (i.e. \textit{of the costructural morphism of a $Y$ in $X\backslash\mathcal{M}$) as $p_f i_f$, with $i_f \in C$, $p_f \in W$. We assume for this that we have}
\msp{page 113}
\textit{a factorization of this type for $X \amalg X \xrightarrow{(\,\mathrm{id}\,,\mathrm{id}\,)} X$,} i.e. \textit{an object $I_X$ and a factorization}
\[
\begin{tikzcd}[column sep=6em]
X \amalg X \arrow[r, hook, "{i_X = (\delta_0^X,\delta_1^X)}"] \arrow[rr, dotted, bend right, "{(\mathrm{id}_X,\mathrm{id}_X)}"'] & I_X \arrow[r, "p_X"] & X\ ,
\end{tikzcd}
\]
\textit{with $i_X \in C$, $p_X \in W$. We then associate to every $f$ the mapping cone} \texttt{[rather cylinder]} \textit{$Z(f)$ as in diagram} (8.6.1) \textit{(where one must read $I_X$ instead of $X \times I$), and we then have $p_f \in W$, $i_f \in C$ as desired. But one must assume, to know that $i_f \in C$, that $X \in \mathcal{M}_C$ ($C$-cofibrant objects),} i.e. \textit{that for every $Y$, $X \longrightarrow Y \amalg X$ is in $C$. (Moreover, we have tacitly assumed $\mathcal{M}$ stable under finite $\varinjlim$.)}
% foot: Version 5 novembre 2010

\printedpage{62}
% head: A. Grothendieck : Dérivateurs, Ch. XV, Édition des Universités de Montpellier II et Paris VII

\textit{The complete diagram is}
\[
\left\{
\begin{array}{c}
\begin{tikzcd}[ampersand replacement=\&,column sep=4em,row sep=3.5em]
 \& X \arrow[r, "f"] \arrow[d, "\mathrm{inj}_0^X"'] \arrow[dr, phantom, "\text{cocart.}"] \& Y \arrow[d] \\
X \arrow[r, "\mathrm{inj}_1^X"] \arrow[dr, dotted, "\delta_1^X"'] \& X \amalg X \arrow[r, "f''"] \arrow[d, "i_X"'] \arrow[dr, phantom, "\text{cocart.}"] \& Y \amalg X \arrow[d, "j"] \\
 \& I_X \arrow[r, "f'"'] \arrow[uu, dotted, bend right=45, "p_X"{pos=.12,swap}] \& Z(f) \arrow[uu, dotted, bend right=45, "p_f"{pos=.25,swap}]
\end{tikzcd}
\\[1ex]
f = p_f i_f, \quad \textit{where} \quad i_f = f' \circ \delta_1^X = j \circ (\underbrace{f'' \circ \mathrm{inj}_1^X}_{\lambda})\ .
\end{array}
\right.
\tag*{(8.7.1)}
\]

\msp{page 114}

We have $i_X \in C$ by hypothesis, hence $j \in C$ (stability of $C$ under cobase change), we have $\lambda : X \longrightarrow Y \amalg X$ in $C$ since $X$ in $\mathcal{M}_C$, hence the composite $i_f = j \circ \lambda$ is in $C$ (stability of $C$ under composition) – we win!

\textbf{Comments.} For the equivalence of $(a)$, $(b)$, $(c)$ in theorem 8.4, $1^\circ$, we thus did not have to call on the developments of Pursuing Stacks, which serve only to establish the equivalence of the three conditions in question with the ``modelizing'' condition $(d)$. To summarize:
\begin{enumerate}[leftmargin=2.5em]
\item[$1^\circ$)] The implications
\[
(a) \ \Longrightarrow\ (\,(b) \ \Longleftrightarrow\ (c) \ \Longleftrightarrow\ (d)\,)
\]
are valid under just the conditions
\[
\mathrm{W}\,(1,\ 2\,\mathrm{bis},\ 3)\ ,
\]
and for the equivalence of $(a)$ with the other conditions, it suffices to have
\[
\mathrm{W}\,(1\,\mathrm{bis},\ 2\,\mathrm{bis},\ 3)
\]
(where W\,(1\,bis) strengthens W\,(1) by the direct factor condition, whereas W\,(2\,bis) weakens W\,(2), requiring only that $W$-aspheric morphisms be in $W$). The conclusions of 8.4, $3^\circ$, concerning the case where $A$ is $W$-aspheric, are also valid under the conditions W\,(1, 2\,bis, 3) alone.
\end{enumerate}

\msp{page 115}

\begin{enumerate}[leftmargin=2.5em]
\item[$2^\circ$)] Corollary 8.5, referring to conditions $(b)$, $(c)$, $(d)$ of $1^\circ$ (to the exclusion of $(a)$, which assumes $W$ stable under direct factors), is valid for a $W$ satisfying
\[
\mathrm{W}\,(1,\ 2\,\mathrm{bis},\ 3,\ \underbrace{7\,\mathrm{bis}}_{\substack{\text{square}\\ \text{cocartesian}}}\,)\ .
\]
% foot: Version 5 novembre 2010
\printedpage{63}
% head: A. Grothendieck : Dérivateurs, Ch. XV, Édition des Universités de Montpellier II et Paris VII
\item[$3^\circ$)] For the validity of 8.4, $2^\circ$, \textit{i.e.} to have a Quillen triple $(W_A,C,F)$, we need (almost) the whole range of conditions on $W$,
\[
\mathrm{W}\,(1\,\mathrm{bis},\ 2\,\mathrm{bis},\ 3,\ \underbrace{\ 6\ }_{\text{limits}},\ \underbrace{7\,\mathrm{bis}}_{\substack{\text{square}\\ \text{cocartesian}}},\ \underbrace{\ 8\ }_{\text{accessibility}})\ ^{\footnotemark[42]}\ ,
\]
\footnotetext[42]{\textbf{N.B.} We should renumber the axioms W\,(\ldots).}
to the exclusion of W\,(4), and being content for W\,(2) with the weakened form W\,(2\,bis) and no more. (Even the connectedness axiom W\,(5) was not used, but I had initially included it, since there is little chance that we will ever have to work with $W$'s that do not satisfy it! It might moreover be, in the presence of the other axioms, or only of W\,(1\,bis, 2\,bis, 3, 7\,bis), that if W\,(5) is \textit{not} satisfied, we have $W = \mathrm{Fl}\,\mathrm{Cat}$ in its entirety, hence $W_A = \mathrm{Fl}\,A^\wedge$!)
\end{enumerate}

\msp{page 116}

\section*{9\quad Revision of the notations, for the axioms of a basic localizer\footnote[43]{\textit{Cf.} XVII \texttt{[rather XVI]} for a definitive version (?).}}

\begin{description}[style=sameline,leftmargin=6.5em]
\item[Loc\,(1)] (Saturation).
\begin{enumerate}[leftmargin=2.5em]
\item[$a$)] Isomorphisms are in $W$.
\item[$b$)] If $f$, $g$ are composable, and two among $f$, $g$, $gf$ are in $W$, the third as well.
\item[$c$)] If $X \underset{g}{\overset{f}{\rightleftarrows}} Y$, if $gf$ and $fg$ are in $W$, then $f$ and $g$ as well.
\end{enumerate}
\end{description}
\noindent Strengthened variant:
\begin{description}[style=sameline,leftmargin=6.5em]
\item[Loc\,(1\,bis)] Loc\,(1) plus
\begin{enumerate}[leftmargin=2.5em]
\item[$d$)] $W$ is stable under direct factors.
\end{enumerate}
\item[Loc\,(2)] (Final object axiom). If $X$ in Cat has a final object, then $X$ is $W$-aspheric, \textit{i.e.} $X \longrightarrow e$ is in $W$.
\item[Loc\,(3)] (Localization axiom, weak form). Let $f : X \longrightarrow Y$ be such that for every $y \in Y$, $X/y$ is $W$-aspheric (\textit{i.e.} the morphism $f/y : X/y \longrightarrow Y/y$ deduced from $f$ by base change $Y/y \longrightarrow Y$ is in $W$). Then $f \in W$.
\end{description}
\noindent Strengthened variant of Loc\,(3):
\begin{description}[style=sameline,leftmargin=6.5em]
\item[Loc\,(3\,bis)] (Localization axiom, strong form). Let $f : X \longrightarrow Y$ be a morphism in $\mathrm{Cat}/S$, such that $\forall\ s \in S$, $f/s : X/s \longrightarrow Y/s$ (deduced from $f$ by base change $S/s \longrightarrow S$) is in $W$. Then $f \in W$. (\textbf{N.B.} Unless I am mistaken, it suffices to require it if $X$, $Y$ are split Cat-fibered over $S$, and $f$ $S$-cartesian compatible with the splitting – assuming Loc\,(1, 2)). % [sic] il suffit le l'exiger -> il suffit de l'exiger, "le" for "de"
\end{description}
% foot: Version 5 novembre 2010